\section{System architecture}
\label{sec:architecture}

\subsection{Four layers}

\ttlibname is organised as four layers with a strict one-way dependency
(Figure~\ref{fig:architecture}). Each layer is useful on its own: a researcher testing a new
feedback rule works in layer~1, someone adding a model variant works in layer~3, and an
applied user only ever sees layer~4.

\begin{figure}[htbp]
\centering
\resizebox{\linewidth}{!}{% Figure: the four layers of torchtsetlin and what each one owns.
\begin{tikzpicture}[font=\sffamily\small, node distance=0pt]

  \def\lw{14.6}          % layer width
  \def\lx{0}             % layer left edge

  % ---------------------------------------------------------------- layer 4 (top)
  \begin{scope}[shift={(0,0)}]
    \draw[ttlayer] (\lx,-1.75) rectangle (\lw,0.72);
    \node[tttitle, anchor=north west, align=left] at (\lx+0.18,0.62)
      {\textbf{4\ \ Workflow}\ \ \textcolor{ttgray}{\code{data/}\ \ \code{train/}\ \
       \code{metrics.py}\ \ \code{interpret.py}\ \ \code{viz.py}}};
    \node[ttnode, anchor=north west, text width=2.85cm] at (\lx+0.2,0.12)
      {\textbf{data}\\[1pt]\scriptsize encoders, synthetic\\\scriptsize + torchvision sets};
    \node[ttnode, anchor=north west, text width=2.85cm] at (\lx+3.15,0.12)
      {\textbf{train}\\[1pt]\scriptsize Trainer, callbacks,\\\scriptsize History};
    \node[ttnode, anchor=north west, text width=2.85cm] at (\lx+6.1,0.12)
      {\textbf{metrics}\\[1pt]\scriptsize accuracy, calibration,\\\scriptsize trustworthiness};
    \node[ttnode, anchor=north west, text width=2.85cm] at (\lx+9.05,0.12)
      {\textbf{interpret}\\[1pt]\scriptsize rules, importance,\\\scriptsize explanations};
    \node[ttnode, anchor=north west, text width=2.4cm] at (\lx+12.0,0.12)
      {\textbf{viz}\\[1pt]\scriptsize clause patches,\\\scriptsize heat maps};
  \end{scope}

  % ---------------------------------------------------------------- layer 3
  \begin{scope}[shift={(0,-2.55)}]
    \draw[ttlayer] (\lx,-2.55) rectangle (\lw,0.72);
    \node[tttitle, anchor=north west, align=left] at (\lx+0.18,0.62)
      {\textbf{3\ \ Concrete models}\ \ \textcolor{ttgray}{\code{models/\{classifier,coalesced,regression,conv\}.py}}%
       \ \ \textemdash\ supply an \emph{output layer} and a \emph{feedback policy}};
    \node[ttnode dark, anchor=north west, text width=3.1cm, minimum height=1.3cm] at (\lx+0.2,0.1)
      {TsetlinMachine\\[1pt]\scriptsize class-owned clauses,\\\scriptsize $\pm$ polarity, weights};
    \node[ttnode dark, anchor=north west, text width=3.1cm, minimum height=1.3cm] at (\lx+3.5,0.1)
      {CoalescedTM\\[1pt]\scriptsize shared pool,\\\scriptsize signed weight / output};
    \node[ttnode dark, anchor=north west, text width=3.1cm, minimum height=1.3cm] at (\lx+6.8,0.1)
      {RegressionTM\\[1pt]\scriptsize positive clauses,\\\scriptsize $v\in[0,T]$ to target};
    \node[ttnode, anchor=north west, text width=3.9cm, minimum height=1.3cm, fill=ttamber!22, draw=ttamber]
      at (\lx+10.35,0.1)
      {\textbf{\_ConvMixin}\\[1pt]\scriptsize patch encode, OR over patches,\\
       \scriptsize random matching patch per event};
    \node[ttlabel, anchor=north west, text width=13.9cm] at (\lx+0.2,-1.60)
      {\textcolor{ttorange}{$\Rightarrow$} mixing gives
       \code{ConvTsetlinMachine}, \code{ConvCoalescedTsetlinMachine},
       \code{ConvRegressionTsetlinMachine}, \code{Conv1dTsetlinMachine} --- one mix-in,
       four convolutional models, no duplicated learning code.};
  \end{scope}

  % ---------------------------------------------------------------- layer 2
  \begin{scope}[shift={(0,-6.05)}]
    \draw[ttlayer] (\lx,-1.75) rectangle (\lw,0.72);
    \node[tttitle, anchor=north west, align=left] at (\lx+0.18,0.62)
      {\textbf{2\ \ Stateful core}\ \ \textcolor{ttgray}{\code{models/base.py :: TsetlinMachineBase(nn.Module)}}};
    \node[ttnode, anchor=north west, text width=3.35cm] at (\lx+0.2,0.1)
      {\textbf{state buffers}\\[1pt]\scriptsize \code{ta\_state} $(C,2F)$ int\\
       \scriptsize \code{include}, \code{include\_count}};
    \node[ttnode, anchor=north west, text width=3.35cm] at (\lx+3.75,0.1)
      {\textbf{learning loop}\\[1pt]\scriptsize coerce, lazy init, chunk,\\
       \scriptsize accumulate, commit};
    \node[ttnode, anchor=north west, text width=3.35cm] at (\lx+7.3,0.1)
      {\textbf{regularisation}\\[1pt]\scriptsize drop-clause,\\\scriptsize drop-literal masks};
    \node[ttnode, anchor=north west, text width=3.35cm] at (\lx+10.85,0.1)
      {\textbf{introspection}\\[1pt]\scriptsize clause expressions,\\\scriptsize state summary};
  \end{scope}

  % ---------------------------------------------------------------- layer 1 (bottom)
  \begin{scope}[shift={(0,-8.30)}]
    \draw[ttlayer] (\lx,-1.75) rectangle (\lw,0.72);
    \node[tttitle, anchor=north west, align=left] at (\lx+0.18,0.62)
      {\textbf{1\ \ Stateless tensor mathematics}\ \ \textcolor{ttgray}{\code{functional.py}, \code{utils.py}}%
       \ \ \textemdash\ no state, no autograd, device-agnostic};
    \node[ttnode plain, anchor=north west, text width=2.6cm] at (\lx+0.2,0.1)
      {\code{to\_literals}\\\scriptsize $x \mapsto [x,\lnot x]$};
    \node[ttnode plain, anchor=north west, text width=2.6cm] at (\lx+2.95,0.1)
      {\code{clause\_outputs}\\\scriptsize one matmul};
    \node[ttnode plain, anchor=north west, text width=2.6cm] at (\lx+5.7,0.1)
      {\code{vote\_sums}\\\scriptsize $c\,W$, clipped};
    \node[ttnode plain, anchor=north west, text width=2.85cm] at (\lx+8.45,0.1)
      {\code{type\_i/ii\_counts}\\\scriptsize event counting};
    \node[ttnode plain, anchor=north west, text width=2.85cm] at (\lx+11.45,0.1)
      {\code{apply\_feedback}\\\scriptsize binomial transitions};
  \end{scope}

  % ---------------------------------------------------------------- dependency arrow
  \draw[ttarrow, line width=1.1pt] (\lw+0.42,-10.0) -- (\lw+0.42,0.55);
  \node[ttlabel, rotate=90, anchor=south, text=ttorange] at (\lw+0.28,-4.8) {depends on};

  % PyTorch substrate
  \node[ttnode slate, anchor=north west, text width=\lw cm, align=center] at (\lx,-10.5)
    {\textbf{PyTorch} \textemdash\ tensors, buffers, devices, \code{state\_dict}, \code{DataLoader}, \code{torch.binomial}};
  \draw[ttarrow, draw=ttslate] (\lw/2,-10.47) -- (\lw/2,-10.25);
\end{tikzpicture}
}
\caption{The four layers. Dependencies point downwards only: the learning mathematics does
  not know what a model is, and a model does not know what a trainer is.}
\label{fig:architecture}
\end{figure}

\paragraph{Layer 1 --- \code{functional.py}.}
Stateless tensor functions: literal expansion, clause evaluation, vote sums, feedback
probabilities, negative-class sampling, Type~I and Type~II event counting, and the automaton
transition itself. No state, no autograd, no device assumptions. All learning mathematics
lives here, which means it can be unit-tested against a brute-force reference implementation
and reused by code that does not use the model classes at all.

\paragraph{Layer 2 --- \clsname{TsetlinMachineBase}.}
An \code{nn.Module} that owns the automata and runs the generic learning loop: input
coercion, lazy shape initialisation, chunking, dropout masks, feedback accumulation and the
commit. It knows nothing about classes, labels or images.

\paragraph{Layer 3 --- concrete models.}
Each variant supplies only an \emph{output layer} (how clause outputs become votes) and a
\emph{feedback policy} (which clauses receive which feedback for a labelled example).
Convolution is orthogonal to both and enters as a mix-in, so
\code{class ConvTsetlinMachine(\_ConvMixin, TsetlinMachine)} is the entire definition of the
convolutional classifier's learning behaviour.

\paragraph{Layer 4 --- workflow.}
Data preparation, training, metrics, interpretation and plotting. These depend on the models
through their public API only, so a user's own model subclass gets the whole workflow layer
for free.

\subsection{The learning state}

The learnable state is a \textbf{buffer, not a parameter}, because a Tsetlin machine has no
gradients:

\begin{center}\small
\begin{tabularx}{0.95\linewidth}{@{}L{3.4cm}L{3.5cm}>{\raggedright\arraybackslash}X@{}}
\ttoprule
\thead{Buffer} & \thead{Shape / dtype} & \thead{Role}\\
\ttmidrule
\code{ta\_state} & $(C, 2F)$ \code{int32} & the automata; literal included iff state $\ge N$\\
\code{include} & $(C, 2F)$ float & cached 0/1 mask, derived, non-persistent\\
\code{include\_count} & $(C,)$ long & literals per clause, derived, non-persistent\\
\code{weights} & $(C,)$ or $(C,K)$ \code{int32} & clause weights (weighted\slash coalesced models)\\
\ttbotrule
\end{tabularx}
\end{center}

Only \code{ta\_state} and \code{weights} are persistent; the two caches are recomputed
whenever the state changes, including after \code{load\_state\_dict}. Storing the derived
mask costs $2CF$ floats and removes a comparison from the inner loop of every forward pass,
which is the single hottest path in the library.

Two consequences are worth calling out. First, \code{model.parameters()} is deliberately
empty and warns when it returns nothing, because silently returning an empty iterator to
someone who wrote \code{optim.SGD(model.parameters())} would be worse than an error. Second,
\code{n\_features=None} enables lazy initialisation in the style of
\code{torch.nn.LazyLinear}: the state is allocated on the first batch, and
\code{\_load\_from\_state\_dict} reconstructs the shape from the checkpoint before loading so
that a lazily built model can be restored without ever having seen data.

\subsection{How an update flows}

\code{update(x, y)} is the analogue of \code{loss.backward(); optimizer.step()}, and it
returns the \emph{pre-update} vote sums so that a training loop can log batch accuracy
without a second forward pass. Figure~\ref{fig:pipeline} shows the flow.

\begin{figure}[htbp]
\centering
\resizebox{\linewidth}{!}{% Figure: the update() pipeline, with the subclass hook points marked.
\begin{tikzpicture}[font=\sffamily\small]

  % ================================================================ row 1: entry
  \node[ttnode dark, anchor=west, text width=2.6cm] (upd) at (0,0)
    {\code{update(x, y)}\\[1pt]\scriptsize one learning step};
  \node[ttnode plain, anchor=west, text width=2.45cm] (prep) at (3.5,0)
    {\code{\_prepare}\\[1pt]\scriptsize coerce input, lazy init};
  \node[ttnode plain, anchor=west, text width=2.75cm] (chunk) at (7.0,0)
    {split the batch into \textbf{chunks}\\[1pt]\scriptsize memory budget};
  \node[ttlabel, anchor=west, text width=3.5cm] at (10.8,0)
    {\textcolor{ttorange}{$\rightarrow$} each chunk runs the row below and adds to the
     accumulator};

  \draw[ttarrow] (2.85,0) -- (3.45,0);
  \draw[ttarrow] (6.35,0) -- (6.95,0);

  % ================================================================ row 2: per chunk
  \begin{scope}[shift={(0,-2.75)}]
    \node[ttnode plain, anchor=west, text width=1.75cm] (enc) at (0.15,0)
      {\code{\_encode}\\[1pt]\scriptsize $x \to$ literals};
    \node[ttnode plain, anchor=west, text width=1.95cm] (ev) at (2.35,0)
      {\code{\_evaluate}\\[1pt]\scriptsize clause outputs};
    \node[ttnode plain, anchor=west, text width=1.55cm] (vo) at (4.75,0)
      {\code{\_votes}\\[1pt]\scriptsize $(B,K)$};
    \node[ttnode, fill=ttamber!25, draw=ttamber, anchor=west, text width=2.35cm] (sel) at (6.75,0)
      {\code{\_select\_feedback}\\[1pt]\scriptsize Type I / II masks};
    \node[ttnode plain, anchor=west, text width=2.35cm] (cnt) at (9.55,0)
      {\code{\_feedback\_counts}\\[1pt]\scriptsize two matmuls};
    \node[ttnode plain, anchor=west, text width=1.9cm] (aw) at (12.35,0)
      {\code{\_accumulate}\\\code{\_weights}};

    \draw[ttarrow] (2.0,0) -- (2.3,0);
    \draw[ttarrow] (4.4,0) -- (4.7,0);
    \draw[ttarrow] (6.4,0) -- (6.7,0);
    \draw[ttarrow] (9.2,0) -- (9.5,0);
    \draw[ttarrow] (12.0,0) -- (12.3,0);

    \node[ttlabel, anchor=west, text=ttdeep] at (0.15,1.05)
      {\textbf{\_accumulate\_chunk} \textemdash\ runs once per chunk; \emph{nothing is written to the automata yet}};
    \begin{scope}[on background layer]
      \draw[ttorange!35, dashed, line width=0.8pt, rounded corners=4pt]
        (-0.2,-0.85) rectangle (14.45,0.8);
    \end{scope}
  \end{scope}

  \draw[ttarrow] (8.0,-0.45) .. controls (8.0,-1.0) and (0.4,-1.1) .. (0.4,-2.3);

  % ================================================================ row 3: accumulator
  \node[ttnode, anchor=west, text width=13.8cm, align=center, fill=ttsand] (acc) at (0.15,-5.05)
    {\textbf{FeedbackAccumulator} \textemdash\ running event counts for the whole update:\\[2pt]
     \code{n\_true} $(C,2F)$ \quad \code{n\_false} $(C,2F)$ \quad \code{n\_ib} $(C)$ \quad
     \code{n2} $(C,2F)$ \quad \code{extra} (weight deltas)};
  \draw[ttarrow] (10.7,-3.65) -- (10.7,-4.6);
  \draw[ttarrow] (13.3,-3.65) -- (13.3,-4.6);

  % ================================================================ row 4: commit
  \begin{scope}[shift={(0,-6.95)}]
    \node[ttnode dark, anchor=west, text width=1.85cm] (com) at (0.15,0) {\code{\_commit}};
    \node[ttnode, anchor=west, text width=2.95cm, align=center] (apf) at (2.5,0)
      {\code{functional.}\\\code{apply\_feedback}\\[1pt]\scriptsize binomial draws, in place};
    \node[ttnode plain, anchor=west, text width=2.35cm] (cw) at (6.4,0)
      {\code{\_commit\_weights}};
    \node[ttnode plain, anchor=west, text width=2.65cm] (ri) at (9.4,0)
      {\code{\_refresh\_include}\\[1pt]\scriptsize rebuild cached masks};
    \node[ttnode dark, anchor=west, text width=2.1cm] (ret) at (12.7,0)
      {\scriptsize returns pre-update votes};
    \draw[ttarrow] (2.15,0) -- (2.45,0);
    \draw[ttarrow] (6.05,0) -- (6.35,0);
    \draw[ttarrow] (9.05,0) -- (9.35,0);
    \draw[ttarrow] (12.35,0) -- (12.65,0);
  \end{scope}
  \draw[ttarrow] (0.7,-5.55) -- (0.7,-6.55);

  % ================================================================ legend
  \node[ttnode, fill=ttamber!25, draw=ttamber, anchor=west, font=\scriptsize, inner sep=4pt]
    (lg1) at (0.15,-8.75) {subclass hook};
  \node[ttlabel, anchor=west, text width=11.6cm] at (2.5,-8.75)
    {A new model implements \code{\_votes}, \code{\_select\_feedback} and
     \code{\_coerce\_targets}; clause weights add
     \code{\_accumulate\_weights}\slash\code{\_commit\_weights}; a different \emph{input
     structure} overrides \code{\_encode}, \code{\_evaluate} and \code{\_feedback\_counts}.
     The loop, the chunking, the dropout masks and the commit are inherited unchanged.};
\end{tikzpicture}
}
\caption{The \code{update()} pipeline. The batch is split into chunks that respect a memory
  budget; each chunk contributes \emph{event counts} to a shared accumulator; the automata
  are written once, at the commit. Amber marks the hook a new model must implement.}
\label{fig:pipeline}
\end{figure}

The separation between \emph{accumulate} and \emph{commit} is the structural core of the
library. Accumulation is pure counting and can be done in any number of chunks, in any order,
without changing the result. The commit is the only place the automata are written, the only
place randomness enters the state, and the only place the derived caches are refreshed. That
split is what makes the batched algorithm expressible, what bounds memory usage, and what
makes the convolutional mix-in possible --- convolution changes what is counted and how, but
not when or how it is committed.

\subsection{Extension points}

A new model is a subclass, not a fork. The required surface is three methods:

\begin{center}\small
\begin{tabularx}{0.95\linewidth}{@{}L{4.6cm}>{\raggedright\arraybackslash}X@{}}
\ttoprule
\thead{Method} & \thead{Contract}\\
\ttmidrule
\code{\_votes(clause\_out)} & $(B,C)$ Boolean clause outputs $\to$ $(B,K)$ vote sums\\
\code{\_select\_feedback(votes, y, clause\_out)} & $\to$ two $(B,C)$ Boolean masks: who gets Type~I, who gets Type~II\\
\code{\_coerce\_targets(y, n, device)} & validate and convert the targets\\
\ttmidrule
\multicolumn{2}{@{}l@{}}{\emph{optional}}\\
\code{\_accumulate\_weights} / \code{\_commit\_weights} & clause weights\\
\code{\_encode} / \code{\_evaluate} / \code{\_feedback\_counts} & a different input structure (this is the convolutional mix-in)\\
\code{\_chunk\_elements\_per\_example} & the memory model for that structure\\
\code{rules()} & how the model prints itself\\
\ttbotrule
\end{tabularx}
\end{center}

\begin{ttkey}
Learning mathematics is stateless and lives in one file; state and the loop live in one base
class; variants supply an output layer and a feedback policy. The entire convolutional family
is one mix-in over that structure --- roughly 360 lines, none of which re-implement learning.
\end{ttkey}
